\documentclass[11pt]{article}
\usepackage[T1]{fontenc}
\usepackage{lmodern}
\usepackage{amsmath,amssymb,amsthm}
\usepackage{enumitem}
\usepackage{graphicx}
\usepackage[margin=1in]{geometry}
\usepackage[colorlinks=true,linkcolor=blue,citecolor=blue,urlcolor=blue]{hyperref}

\newtheorem{theorem}{Theorem}
\newtheorem{lemma}{Lemma}
\newtheorem{remark}{Remark}
\newcommand{\capc}{\operatorname{cap}_2}
\newcommand{\diam}{\operatorname{diam}}
\newcommand{\dv}{\,d v_g}

\title{Smooth Spectral Approximation of Compact Trace Measures\\
for Uniformly Elliptic Energy Densities}
\author{Run \texttt{fa\_banach\_001}, agent \texttt{agent\_lane\_14}}
\date{19 July 2026}

\begin{document}
\maketitle

\section*{Status}

Candidate partial result, likely valid. This packet answers the smooth-measure
approximation question in Vinokurov's Remark~1.3 when the energy density is
continuous and uniformly positive. It does not answer the rough or degenerate
weighted case.

\section*{Source Question}

Vinokurov asks in Remark~1.3 of arXiv:2601.17896v3, PDF page~4, whether the
eigenvalues of an arbitrary measure inducing a compact map
\(H^1(M,\alpha)\to L^2(M,\mu)\) can be approximated, or at least bounded from
above, using smooth measures~\cite{Vinokurov}. The source passage is shown
below.

\begin{center}
\includegraphics[width=\linewidth]{figures/open_problem_crop.png}
\end{center}

The packet proves convergence, not merely an upper bound, under the additional
hypothesis that \(\alpha\) is continuous and bounded away from zero.

\section*{Definitions}

Let \((M,g)\) be a closed connected Riemannian manifold of dimension
\(m\geq2\). For a positive density \(\alpha\) and a finite positive Radon
measure \(\mu\), write
\[
 a_\alpha(u,v)=\int_M \alpha\langle du,dv\rangle_g\dv,
 \qquad
 b_\mu(u,v)=\int_M \widetilde u\,\widetilde v\,d\mu,
\]
where tildes denote quasi-continuous representatives. The variational
eigenvalues are
\[
 \lambda_k(\alpha,\mu)
 =\inf_{\substack{F\subset C^\infty(M)\\ \dim F=k+1}}
   \sup_{0\ne u\in F}
   \frac{a_\alpha(u,u)}{b_\mu(u,u)},
\]
with the usual restriction that \(F\) remain \((k+1)\)-dimensional in
\(L^2(\mu)\).

For a Borel set \(E\subset M\), let \(\capc(E)\) be its Sobolev
\(2\)-capacity, and define the local capacitary modulus
\[
 \omega_\sigma(r)=
 \sup\left\{\frac{\sigma(E)}{\capc(E)}:
       E\subset M,\ \diam(E)\leq r,\ \capc(E)>0\right\}.
\]
Maz'ya's compactness criterion states that the trace
\(H^1(M)\to L^2(M,\sigma)\) is compact if and only if
\(\omega_\sigma(r)\to0\) as \(r\downarrow0\); see
\cite[Section~11.9.1]{Mazya} and the manifold formulation recalled in
\cite[Section~2.2]{GKL}.

\section*{Proof Intuition}

Weak-star smoothing alone has the wrong one-sided min--max behavior. The key
is to smooth \(\mu\) without losing the quantitative local property that makes
its Sobolev trace compact. We localize \(\mu\), average each piece under small
diffeomorphisms, and use the quasi-invariance of capacity to keep one uniform
vanishing modulus for all smoothed measures. That modulus turns weak
convergence in \(H^1\) into strong convergence in the moving spaces
\(L^2(\mu_n)\). Weak-star convergence then upgrades to operator-norm
convergence of the quadratic mass forms. Uniform approximation of \(\alpha\)
and compact self-adjoint perturbation theory finish the eigenvalue argument.

\section*{Main Result}

\begin{theorem}\label{thm:main}
Let \((M,g)\) be closed and connected, let \(\alpha\in C(M)\) satisfy
\(0<c\leq\alpha\leq C<\infty\), and let \(\mu\ne0\) be a finite positive
Radon measure such that
\[
 H^1(M,\alpha)\longrightarrow L^2(M,\mu)
\]
is compact. Then there are \(\alpha_n,\rho_n\in C^\infty(M)\), with
\(\alpha_n>0\), \(\rho_n>0\), such that, putting
\(\mu_n=\rho_n v_g\),
\begin{align*}
 &\alpha_n\longrightarrow\alpha \quad\hbox{uniformly},\\
 &\mu_n\stackrel{*}{\rightharpoonup}\mu,
 \qquad \mu_n(M)=\mu(M),\\
 &\sup_{\|u\|_{H^1}\leq1}
   \left|b_{\mu_n}(u,u)-b_\mu(u,u)\right|\longrightarrow0.
\end{align*}
Consequently, for every fixed finite variational eigenvalue,
\[
 \lambda_k(\alpha_n,\mu_n)\longrightarrow\lambda_k(\alpha,\mu).
\]
The corresponding normalized eigenvalues converge as well.
\end{theorem}

\section*{Proof}

Uniform ellipticity makes \(H^1(M,\alpha)\) and the ordinary \(H^1(M)\)
the same Hilbert space with equivalent norms. Hence the assumed trace is
compact exactly when the ordinary Sobolev trace is compact. Maz'ya's criterion
therefore gives
\begin{equation}\label{eq:vanish}
 \omega_\mu(r)\longrightarrow0.
\end{equation}

\subsection*{Step 1: smooth while preserving the capacitary modulus}

Choose a finite smooth partition of unity \(\sum_i\psi_i=1\). On a
neighborhood of \(\operatorname{supp}\psi_i\), choose compactly supported
smooth vector fields \(X_{i,1},\ldots,X_{i,m}\) forming a local frame. For
small \(z\in\mathbb R^m\), let \(\Phi_{i,z}\) be the composition of the
time-\(z_j\) flows of these vector fields. The maps \(\Phi_{i,z}\) are global
diffeomorphisms, converge to the identity with \(z\), and have uniformly
bounded bi-Lipschitz constants for small \(z\).

Let \(\eta_t\) be a nonnegative smooth approximate identity supported in a
small ball in \(\mathbb R^m\), and set
\begin{equation}\label{eq:smoothing}
 \nu_t=\sum_i\int_{\mathbb R^m}
       (\Phi_{i,z})_*(\psi_i\mu)\,\eta_t(z)\,dz.
\end{equation}
Because \(z\mapsto\Phi_{i,z}(x)\) is a local submersion on the support of
\(\psi_i\), each averaged measure in \eqref{eq:smoothing} has a smooth
density. Thus \(\nu_t\) has a nonnegative smooth density,
\(\nu_t(M)=\mu(M)\), and \(\nu_t\stackrel{*}{\rightharpoonup}\mu\).

Sobolev capacity is quasi-invariant under a uniformly bi-Lipschitz family of
diffeomorphisms. Thus there are constants \(A,B\), independent of small
\(t,z\), such that for every set \(E\) of sufficiently small diameter,
\begin{align*}
 (\psi_i\mu)(\Phi_{i,z}^{-1}E)
 &\leq \mu(\Phi_{i,z}^{-1}E)\\
 &\leq \omega_\mu(A\diam E)\,
        \capc(\Phi_{i,z}^{-1}E)\\
 &\leq B\,\omega_\mu(A\diam E)\,\capc(E).
\end{align*}
After summing and averaging,
\begin{equation}\label{eq:uniformmodulus}
 \omega_{\nu_t}(r)\leq B'\omega_\mu(Ar).
\end{equation}
To make the density strictly positive while preserving mass, replace
\(\nu_t\) by
\[
 \widehat\nu_t=(1-\varepsilon_t)\nu_t
   +\varepsilon_t\frac{\mu(M)}{v_g(M)}v_g,
 \qquad \varepsilon_t\downarrow0.
\]
The volume measure has a vanishing capacitary modulus, so
\eqref{eq:uniformmodulus} and \eqref{eq:vanish} remain uniform for
\(\widehat\nu_t\). Choose \(t_n\downarrow0\) and put
\(\mu_n=\widehat\nu_{t_n}\).

\subsection*{Step 2: uniform modulus gives collective compactness}

We record the quantitative consequence needed below. Suppose positive measures
\(\sigma_n\) have uniformly bounded masses and
\begin{equation}\label{eq:familymodulus}
 \Omega(r):=\sup_n\omega_{\sigma_n}(r)\longrightarrow0.
\end{equation}
If \(u_n\rightharpoonup0\) in \(H^1(M)\), then
\begin{equation}\label{eq:collective}
 \int_M u_n^2\,d\sigma_n\longrightarrow0.
\end{equation}

Indeed, cover \(M\) by finitely many balls \(B_j\) of radius \(r\), with the
doubled balls having bounded overlap. For every Borel set \(E\),
\[
 (\sigma_n|_{B_j})(E)
 =\sigma_n(E\cap B_j)
 \leq \Omega(2r)\capc(E\cap B_j)
 \leq \Omega(2r)\capc(E).
\]
The capacitary trace inequality therefore gives a trace norm of order
\(\Omega(2r)^{1/2}\) for \(\sigma_n|_{B_j}\). Apply it to a cutoff times
\(u_n-(u_n)_{2B_j}\), and then use the Poincare inequality, to obtain
\[
 \int_{B_j}|u_n-(u_n)_{2B_j}|^2\,d\sigma_n
 \leq C\Omega(2r)\int_{2B_j}|du_n|^2\dv.
\]
For fixed \(r\), weak convergence implies
\((u_n)_{2B_j}\to0\) for every one of the finitely many balls. Summing yields
\[
 \limsup_n\int_Mu_n^2\,d\sigma_n
 \leq C\Omega(2r)\sup_n\|u_n\|_{H^1}^2.
\]
Letting \(r\downarrow0\) proves \eqref{eq:collective}. The same localization
argument, without assuming weak convergence, also gives a uniform bound for
the trace norms \(H^1(M)\to L^2(M,\sigma_n)\).

By \eqref{eq:uniformmodulus}, this lemma applies to the constructed
\(\mu_n\).

\subsection*{Step 3: convergence in Sobolev form norm}

First, for every fixed \(u\in H^1(M)\),
\begin{equation}\label{eq:fixed}
 b_{\mu_n}(u,u)\longrightarrow b_\mu(u,u).
\end{equation}
For smooth \(u\), this is weak-star convergence. For general \(u\), approximate
it in \(H^1\) by smooth functions and use the uniform trace bounds for
\(\mu_n\), together with the bounded trace for \(\mu\).

Suppose the asserted form-norm convergence were false. After passing to a
subsequence, there would be \(u_n\) in the \(H^1\) unit ball and a fixed
\(\delta>0\) such that
\[
 |b_{\mu_n}(u_n,u_n)-b_\mu(u_n,u_n)|\geq\delta.
\]
Take a further subsequence with \(u_n\rightharpoonup u\) in \(H^1\), and set
\(v_n=u_n-u\). Compactness of the limiting trace gives
\(b_\mu(v_n,v_n)\to0\); collective compactness gives
\(b_{\mu_n}(v_n,v_n)\to0\). The mixed terms tend to zero by
Cauchy--Schwarz, while \eqref{eq:fixed} handles the fixed vector \(u\). This
contradicts the choice of \(u_n\), and proves
\begin{equation}\label{eq:formnorm}
 \|b_{\mu_n}-b_\mu\|_{\mathcal B(H^1\times H^1)}\longrightarrow0.
\end{equation}

\subsection*{Step 4: pass to eigenvalues}

Choose positive \(\alpha_n\in C^\infty(M)\) with
\(\alpha_n\to\alpha\) uniformly. Then
\[
 \|a_{\alpha_n}-a_\alpha\|_{\mathcal B(H^1\times H^1)}
 \leq \|\alpha_n-\alpha\|_\infty\longrightarrow0.
\]
The form
\(h=a_\alpha+b_\mu\) is an equivalent Hilbert norm on \(H^1(M)\): this is
the standard Poincare inequality with the nonzero continuous functional
\(u\mapsto\int_Mu\,d\mu\). Put
\(h_n=a_{\alpha_n}+b_{\mu_n}\). Relative to \(h\), let \(G_n\) and \(D_n\)
represent \(h_n\) and \(b_{\mu_n}\), respectively. Equations
\eqref{eq:formnorm} and the uniform convergence of \(\alpha_n\) give
\[
 G_n\longrightarrow I,
 \qquad D_n\longrightarrow D
 \quad\hbox{in operator norm},
\]
where \(D\) represents \(b_\mu\). Hence the compact positive self-adjoint
operators
\[
 S_n=G_n^{-1/2}D_nG_n^{-1/2}
\]
converge in norm to \(D\). Their ordered eigenvalues therefore converge. If
\(\theta_{k,n}>0\) is an eigenvalue of \(S_n\), the generalized eigenvalue
relation is
\[
 b_{\mu_n}(u,v)=\theta_{k,n}h_n(u,v),
 \qquad
 \lambda_k(\alpha_n,\mu_n)=\theta_{k,n}^{-1}-1.
\]
The same identity holds at the limit, so every fixed finite variational
eigenvalue converges. Since the masses are equal and the relevant norms of
\(\alpha_n\) converge, normalized eigenvalues converge as well. This proves
Theorem~\ref{thm:main}. \qed

\section*{Verification Notes}

The separate file \texttt{verification.md} records the adversarial check. The
main points are:
\begin{enumerate}[leftmargin=*]
\item ordinary weak-star smoothing alone is insufficient, and the proof does
not use it as though it implied spectral convergence;
\item the uniform capacity estimate follows before any spectral argument and
is stable under adding a vanishing positive volume component;
\item collective compactness is proved for moving measures, which is the step
needed to upgrade fixed-vector convergence to operator-norm convergence;
\item the generalized eigenproblem is conjugated to self-adjoint compact
operators on one fixed Hilbert space before perturbation theory is invoked.
\end{enumerate}

\section*{Limitations and Novelty Check}

The source question allows rough or degenerate \(\alpha\). The present proof
uses uniform ellipticity twice: to invoke the ordinary Sobolev capacity
criterion and to approximate the energy form in operator norm. For a rough
weight, local diffeomorphisms need not preserve weighted capacity; for a
degenerate weight, even the underlying Sobolev space can change under
smoothing. Those cases remain open.

Before promotion, the run's four lightweight indexes were searched for the
arXiv id and the exact spectral-approximation terms. A bounded web search used
the exact question phrase, the source title/id, and combinations of compact
Sobolev trace, capacity, mollification, smooth densities, and eigenvalue
convergence. It found the source paper and known continuity results, but no
later explicit answer or this exact approximation theorem. The main known
continuity reference, Girouard--Karpukhin--Lagace, proves convergence under
Sobolev-dual convergence and supplies the admissible-measure framework; it
does not state the arbitrary compact-trace smoothing conclusion recorded
here~\cite{GKL}.

\section*{Human Review Recommendation}

Review as a substantial partial result. The highest-value checks are the
smooth localized-flow construction, the uniform capacity comparison in
\eqref{eq:uniformmodulus}, and the localized proof of collective compactness.
Do not promote the packet to a full answer without an argument for rough or
degenerate \(\alpha\).

\begin{thebibliography}{9}

\bibitem{Vinokurov}
D. Vinokurov,
\emph{Eigenvalue optimization in higher dimensions and \(p\)-harmonic maps},
arXiv:2601.17896v3 (2026).

\bibitem{GKL}
A. Girouard, M. Karpukhin, and J. Lagace,
\emph{Continuity of eigenvalues and shape optimisation for Laplace and
Steklov problems},
Geom. Funct. Anal. \textbf{31} (2021), 513--561;
arXiv:2004.10784.

\bibitem{Mazya}
V. G. Maz'ya,
\emph{Sobolev Spaces with Applications to Elliptic Partial Differential
Equations},
Grundlehren der mathematischen Wissenschaften 342, augmented edition,
Springer, 2011.

\end{thebibliography}

\end{document}
